\section{An explicit finite threshold}\label{app:explicit-bound}

The main proof needs only the qualitative limits in \Cref{B:prop:R}. Here we retain the explicit constants from the quantitative argument and obtain a sufficient finite order.

\begin{theorem}[Uniform finite-scale non-SoS test]\label{thm:finite-scale}
Let $k,\epsilon,m$ satisfy \eqref{B:eq:dyadic-scale}. If $\calF_m$ is SoS,
then every positive semidefinite normal-form Gram representation \eqref{eq:general-gram-1} gives a
nonnegative scalar $\Phi_m$ satisfying
\begin{equation}\label{eq:finite-scale}
          0\le\Phi_m
          <-2^{42}+2^{90}m^{-1/2}+2^{80}m^{-1/262144}.
\end{equation}
All constants are independent of the Gram corrections and of the rank
of the representation.
\end{theorem}
\begin{proof}

Use the fixed data of \eqref{C:eq:witness}, and put
$\mathsf x_\ell=(\xi_\ell,\eta_\ell)
=(ie^{-\epsilon a_\ell},-ie^{-\epsilon\bar a_\ell})$.
Then $\eta_\ell=\bar\xi_\ell$.
Define
\begin{equation}\label{eq:evaluated-gram}
 \mathcal G_m=
 \left[
 \left(M+\frac{P_++P_-}{2}\right)(\overline{\mathsf x_\ell},\mathsf x_h)
 \right]_{\ell,h\in I},
 \qquad \Phi_m=\epsilon^4d^*\mathcal G_md.
\end{equation}
Each summand is an evaluated positive semidefinite kernel on the same
list of points. Taking the direct sum of the pure and mixed feature
spaces proves $\mathcal G_m\succeq0$, hence $\Phi_m\ge0$.

Let $\mathbf L_m$ and $\mathbf R_m$ be the evaluation matrices of
$L_m$ and $\calR P_-$ at these points. Realignment gives
\begin{equation}\label{eq:realigned-diagonal}
 (\calR P_-)(\overline{\mathsf x_\ell},\mathsf x_h)
 =P_-(\bar\xi_\ell,\xi_h,\xi_\ell,\bar\xi_h)
 =P_-[\bar\xi_\ell,\xi_h]=P_-[\xi_\ell,\bar\xi_h].
\end{equation}
The last equality uses real symmetry of $P_-$. For $\ell=(r,j,\sigma)$ and $h=(r',j',\sigma')$, take
\[
 \alpha=\frac{r+r'}2+i\frac{\sigma j-\sigma'j'}2,
 \qquad
 \beta=\frac{\sigma j+\sigma'j'}2-i\frac{r-r'}2.
\]
Then $\alpha+i\beta=a_\ell$ and $\alpha-i\beta=\bar a_h$.
These parameters satisfy \eqref{B:eq:speed-region}, so
\Cref{B:lem:quant-continuation} applies to every entry of
\eqref{eq:realigned-diagonal}. Consequently,
\begin{equation}\label{eq:realigned-error}
 |\epsilon^4d^*\mathbf R_md|
 \le\norm d_1^2\,2^{20}\epsilon^{2^{-17}}
 <2^{80}\epsilon^{2^{-17}}.
\end{equation}

It remains to compare the finite baseline with $\mathbf L_\star$. At the test
points, each factor of $A$ and $C$ has the form $1-e^{-\epsilon\chi}$,
where $\Re\chi\ge2$ and $|\chi|<2^{10}$.
The integral identity
\[
 \frac{1-e^{-\epsilon\chi}}{\epsilon\chi}
 =\int_0^1e^{-\epsilon\chi u}\,du
\]
and $|e^{-\epsilon\chi u}-1|\le\epsilon|\chi|u$ give a relative
error at most $2^{10}\epsilon$ for each factor. Thus
$A/\epsilon^2=A_0(1+\delta_A)$ and
$C/\epsilon^2=C_0(1+\delta_C)$, where
$|\delta_A|,|\delta_C|\le2^{12}\epsilon\le1/16$.

For $|\delta|\le1/4$ we use
$|(1+\delta)^{-1}-1|\le2|\delta|$ and
$|(1+\delta)^{-2}-1|\le6|\delta|$.
At our points $A_0,C_0\ge4$. Each factor of $\Delta$ is
$1+e^{-\epsilon\chi}$, and $|\Im(\epsilon\chi)|\le512\epsilon<1$,
so the exponential has positive real part and $|\Delta|\ge1$.
Writing $\delta=\max(|\delta_A|,|\delta_C|)$, the product of the two
inverse factors has relative error at most $5\delta$. The three reciprocal
terms in \eqref{B:eq:generating} therefore contribute at most
$(3/2+3/4+5/8)\delta<3\delta\le3\cdot2^{12}\epsilon$; the remaining
$\Delta$ term contributes at most $4\epsilon^4$. In particular,
\[
 |\epsilon^4 L_\infty(\overline{\mathsf x_\ell},\mathsf x_h)-\mathbf L_\star(\ell,h)|
 \le2^{20}\epsilon.
\]
All four node moduli are at most $e^{-\epsilon}$, so
\Cref{B:lem:quant-tails} and \eqref{B:eq:dyadic-tails} add at most
$\epsilon$ to this scaled error. We obtain
\begin{equation}\label{eq:finite-tangent}
 |\epsilon^4 L_m(\overline{\mathsf x_\ell},\mathsf x_h)-\mathbf L_\star(\ell,h)|
 \le2^{30}\epsilon.
\end{equation}

Pairing \eqref{eq:finite-decomposition} with $d$, and using
\eqref{C:eq:witness-bounds}, \eqref{eq:realigned-error} and
\eqref{eq:finite-tangent}, gives
\begin{align*}
 0\le\Phi_m
 &=\epsilon^4d^*\mathbf L_md-\epsilon^4d^*\mathbf R_md\\
 &<-2^{42}+2^{90}\epsilon+2^{80}\epsilon^{2^{-17}}\\
 &=-2^{42}+2^{90}m^{-1/2}+2^{80}m^{-1/262144}.
\end{align*}
Here both entrywise errors are multiplied by
$\norm d_1^2<2^{60}$, and $\epsilon=m^{-1/2}$.
This proves \eqref{eq:finite-scale}.
\end{proof}

\begin{proposition}[Explicit sufficient order]\label{prop:explicit-threshold}
Let
\[
 k=38\cdot2^{17}+1=\ScaleExponent,\qquad
 m=2^{2k}=\CornerDepth.
\]
Then $\mathcal F_m$ is not SoS. Consequently, $F_N$ is not SoS for
every
\[
                       N\ge2m=\SufficientOrder.
\]
\end{proposition}
\begin{proof}
Convexity of $x\mapsto2^{-x}$ gives
$2^{-x}\le1-x/2$ for $0\le x\le1$. The upper bound in
\eqref{eq:finite-scale} is therefore at most
\[
 -2^{42}+2^{90-k}+2^{42}2^{-2^{-17}}
 \le-2^{24}+2^{90-k}<0.
\]
This contradicts $\Phi_m\ge0$, so $\calF_m$ is not SoS. By
\Cref{lem:corner}, it is a linear restriction of $F_N$ for every
$N\ge2m=\SufficientOrder$, and any SoS representation of such an
$F_N$ would restrict to one of $\calF_m$.
\end{proof}

The size of this threshold comes from the slow term
$2^{80}m^{-1/262144}$ in \eqref{eq:finite-scale}. The witness itself is
fixed; the large order is needed only to make the uniform continuation
error smaller than its negative margin. This bound is sufficient and is
not intended to identify the first non-SoS order.
